\chapter[Dependent inference and causality]{Dependent observations, selective inference, and causality}
\label{ch:dependent-inference}
Time bins from one animal are typically correlated.
Repeating a test at many channels, lags and frequency bands
also creates selection effects. A valid mathematical result
must state which dependence and selection operations it covers.

\section{Variance of a time-series average}
For a zero-mean weakly stationary time series with
$\gamma(k)=\Cov(X_t,X_{t+k})$, direct expansion gives
\[
\Var(\bar X_T)=\frac{1}{T}
\left[\gamma(0)+2\sum_{k=1}^{T-1}
\left(1-\frac{k}{T}\right)\gamma(k)\right].
\]
This \emph{finite-sample identity} needs no central limit theorem.
When $\sum_k|\gamma(k)|<\infty$, the asymptotic variance is
$[\gamma(0)+2\sum_{k\ge1}\gamma(k)]/T$.
For normalized autocorrelation $\rho(k)$, an approximate effective
sample size is
$T_{\mathrm{eff}}=T/[1+2\sum_{k\ge1}\rho(k)]$,
provided that series is meaningful. Negative correlations can make
this heuristic exceed $T$; it must not be blindly clipped or treated
as an exact count of independent animals.

\section{HAC and resampling under dependence}
A heteroskedasticity-and-autocorrelation-consistent (HAC)
variance estimator uses weighted autocovariances, for example
\[
\widehat\Omega=\widehat\gamma(0)
 +2\sum_{k=1}^{L}\left(1-\frac{k}{L+1}\right)
 \widehat\gamma(k).
\]
This is a Bartlett-window Newey--West form in a scalar,
mean-estimation setting. Consistency requires hypotheses on
stationarity, weak dependence, moments and bandwidth growth;
finite samples and nonstationary trials require scrutiny.
Block bootstrap preserves within-block temporal structure
but is only justified under suitable dependence assumptions
and a suitable block length. Randomly permuting every LFP
sample is generally \emph{not} an exchangeable null.

\section{Multiple comparisons and post-selection}
Holm's step-down procedure controls family-wise error under
valid marginal $p$-values without an independence assumption.
Benjamini--Yekutieli controls FDR under arbitrary dependence
using the harmonic correction $c_m=\sum_{j=1}^m1/j$
in its standard formulation; it may be conservative.
A maximum-statistic permutation test controls a family of tests
if the randomization acts on the units exchangeable under the
joint null. Cluster-permutation procedures use adjacency
and a cluster-level statistic; rejection does not locate the
exact significant onset, offset, or single electrode.
After selecting the largest spectral peak and testing it
on the same data, ordinary unadjusted $p$-values generally
fail to reflect that selection. Sample splitting, independent
replication or mathematically justified selective inference
are possible remedies.

\section{Observational dependence versus intervention}
A directed acyclic graph (DAG) encodes a collection of
conditional-independence relations under the Markov property.
Paths are blocked or unblocked according to \emph{d-separation}.
If $Z$ is an appropriate backdoor adjustment set for a proposed
causal effect $X\to Y$, under causal sufficiency assumptions
relative to the DAG and positivity,
\[
p(y\mid\operatorname{do}(x))
=\int p(y\mid x,z)p(z)\,dz.
\]
The formula fails for invalid adjustment sets, unmeasured
confounding, or violated structural assumptions. In particular,
a lagged correlation does not supply the DAG.
Faithfulness is an additional assumption (and can fail by
parameter cancellation); the common-cause principle motivates
but does not uniquely identify latent confounders.

\section{Transfer entropy and Gaussian Granger}
Transfer entropy from $X$ to $Y$, conditional on selected
histories $Y^-$ and $Z^-$, is a conditional mutual information:
$TE=I(Y_t;X^-\mid Y^-,Z^-)$.
For jointly Gaussian variables, with unrestricted correct
linear conditional predictors and covariance determinants
specified consistently,
\[
TE_{X\to Y\mid Z}=\tfrac12
\log\frac{\det\Sigma_{\mathrm{reduced}}}
                {\det\Sigma_{\mathrm{full}}}
\quad\text{(nats)}.
\]
This Gaussian equivalence is established in \cite{barnett2009}.
A common multivariate Granger measure is the full log ratio;
therefore it equals $2TE$ in this Gaussian setting.
This is a theorem about \emph{predictive distributions}, not an
identification theorem for causal interventions.
Spectral Granger additionally requires consistent VAR transfer
functions or spectral factorization (e.g., Wilson methods) and
model residual checks. The present
\path{granger_causality.m} computes a temporal
log-variance ratio, not spectral Granger; the existing
\path{transfer_entropy.m} is a first-order,
discrete plug-in conditional-information estimator.
\begin{caution}The two MATLAB primitives are exploratory and do not by themselves establish causal identifiability.\end{caution}
% --- expanded edition (October 2026) ---
\section{Holm's procedure controls the family-wise error}
\begin{proposition}
Let $p_{(1)}\le\cdots\le p_{(m)}$ be ordered $p$-values with hypotheses
$H_{(i)}$. Reject $H_{(1)},\ldots,H_{(k-1)}$ where $k$ is the first index
with $p_{(k)}>\alpha/(m-k+1)$. If every null $p$-value is valid
($\Prob(p\le u)\le u$), then $\Prob(\text{any true null rejected})\le\alpha$,
whatever the dependence.
\end{proposition}
\begin{proof}
Let $I_0$ be the set of $m_0$ true nulls and $E$ the event that every
$p_i$ with $i\in I_0$ exceeds $\alpha/m_0$. By the union bound,
$\Prob(E^c)\le m_0\cdot\alpha/m_0=\alpha$. On $E$, consider the first true
null in the ordered list, at position $j$. At most $m-m_0$ hypotheses
precede it, so $j\le m-m_0+1$ and its threshold is
$\alpha/(m-j+1)\le\alpha/m_0$, which its $p$-value exceeds. The procedure
stops there or earlier and rejects no true null.
\end{proof}

\section{d-separation by example}
Three elementary configurations determine whether a path between $X$
and $Y$ transmits association:
\begin{itemize}[leftmargin=2em]
\item a \emph{chain} $X\to M\to Y$ and a \emph{fork} $X\leftarrow C\to Y$
are open unless the middle node is conditioned on;
\item a \emph{collider} $X\to K\leftarrow Y$ is closed unless $K$ or one of
its descendants is conditioned on.
\end{itemize}
A path is blocked by a conditioning set $Z$ if it contains a non-collider
in $Z$ or a collider with neither it nor any descendant in $Z$.
\begin{example}[Conditioning on a detected event]
Let $X$ and $Y$ be the activities of two independent neurons, and let a
``population event'' $K$ be detected when the summed activity crosses a
threshold. Analysing only event-triggered epochs conditions on $K$, a
collider: within events, the two neurons become negatively associated
(if one is quiet, the other must have been active to trigger the event).
Event-triggered correlations therefore cannot be compared with
unconditional ones without modelling the selection.
\end{example}
\begin{example}[Behaviour as a common cause]
If running speed $C$ modulates both an LFP band $X$ and a unit's rate $Y$,
the fork $X\leftarrow C\to Y$ produces spike--field correlations.
Adjusting for $C$ (the backdoor formula of this chapter) blocks the
path. If speed is measured noisily, adjustment is incomplete, and
residual association persists: measurement error in confounders is a
common reason why ``controlling for behaviour'' fails.
\end{example}

\section{Spectral factorization and Wilson's method}
Assume a zero-mean, finite-variance, purely nondeterministic stationary
process with an absolutely continuous spectral density matrix $S(f)$.
For a full-rank factorization require $S(f)$ Hermitian positive definite
almost everywhere and $\int_{-1/2}^{1/2}\log\det S(f)\,df>-\infty$.
Under these conditions the normalized causal outer factor gives
\[
S(f)=H(f)\Sigma H(f)^\ast,\qquad
H(f)=I+\sum_{k\ge1}H_ke^{-2\pi ifk},\qquad \Sigma\succ0.
\]
The normalization $H_0=I$ identifies the innovation covariance.
``Minimum phase'' refers here to the outer causal factor; absolute
summability of both the filter and its inverse requires stronger regularity
than stationarity alone (for example a stable invertible rational model).
Singular spectra and deterministic spectral lines require separate treatment.

For a scalar positive spectrum with integrable logarithm, in units of
cycles per sample, the Kolmogorov--Szeg\H{o} formula is
$\log\sigma^2=\int_{-1/2}^{1/2}\log S(f)\,df$.
For the multivariate normalized factor it determines
$\log\det\Sigma=\int_{-1/2}^{1/2}\log\det S(f)\,df$, not each entry
of $\Sigma$. The scalar AR(1) calculation below illustrates this identity;
it is not a proof of the general matrix factorization theorem.

Wilson's iterative factorization targets the causal outer factor on a
frequency grid (Wilson, 1972, DOI: \texttt{10.1137/0123044}). In a numerical
implementation check Hermitian symmetry and positive eigenvalues, specify
initialization and normalization, and report a reconstruction residual and
convergence tolerance. Near-singular estimated spectra can require
regularization, which changes the spectrum being factorized. A finite grid
and finite iterations do not by themselves certify minimum phase or
convergence. The factor may feed a spectral Granger calculation, subject to
the assumptions and uncertainty of the spectral estimate. This manuscript
does not claim a verified Wilson implementation in the current arsenal.

\begin{proposition}[Kolmogorov--Szeg\H{o} for AR(1)]
For $X_t=aX_{t-1}+\varepsilon_t$ with $|a|<1$ and innovation variance
$\sigma^2$, $\int_{-1/2}^{1/2}\log S(f)\,df=\log\sigma^2$.
\end{proposition}
\begin{proof}
$S(f)=\sigma^2/|1-ae^{-2\pi if}|^2$, so the integral is
$\log\sigma^2-2\Re\int\log(1-ae^{-2\pi if})\,df$. Expanding
$\log(1-az)=-\sum_{k\ge1}a^kz^k/k$ and integrating term by term over the
unit circle, every term vanishes.
\end{proof}

\section*{Exercises for this chapter}
\begin{exercise}
Apply Holm and Bonferroni to five $p$-values of your choice and find a
case where Holm rejects more hypotheses.
\end{exercise}
\begin{exercise}
For the DAG $C\to X\to Y$, $C\to Y$, show that $\{C\}$ is a valid
backdoor adjustment set and that $\varnothing$ is not.
\end{exercise}
\begin{exercise}
Simulate two independent Gaussian sequences, select samples where their
sum exceeds 2, and compute their correlation within the selection.
\end{exercise}
\begin{exercise}
Compute the effective sample size of an AR(1) series with $a=0.8$ and
$T=500$.
\end{exercise}
